\documentclass[11pt,a4paper]{article}
\usepackage[margin=1in]{geometry}
\usepackage{amsmath,amssymb}
\usepackage{booktabs}
\usepackage{graphicx}
\usepackage{hyperref}
\usepackage{siunitx}
\sisetup{separate-uncertainty = true}

\title{\textbf{Laboratory Experiment Report:}\\[0.5em] \large Characterization of Nonlinear Dynamic Systems}
\author{
  \textbf{Student Name} \quad (ID: 12345678) \\
  \textit{Department of Electrical and Computer Engineering} \\
  \textit{Lab Partner: Partner Name} \\
  \texttt{student@university.edu}
}
\date{Date of Experiment: \today \quad $|$ \quad Course: ECE-301}

\begin{document}
\maketitle

\begin{abstract}
This laboratory experiment investigates the transient response and frequency behavior of a second-order RLC resonator. Experimental damping ratios and natural frequencies were measured across 10 trials. The empirical resonant frequency was determined to be $\qty{15.42 \pm 0.18}{\kilo\hertz}$, matching theoretical calculations within a relative error margin of $1.8\%$.
\end{abstract}

\section{Objectives}
\begin{itemize}
    \item Measure the step response and transient overshoot of a second-order analog filter.
    \item Calculate empirical damping ratios ($\zeta$) and natural frequencies ($\omega_n$).
    \item Quantify measurement uncertainties and verify theoretical models.
\end{itemize}

\section{Theoretical Background}
The governing differential equation for the series resonator is:
\begin{equation}
L \frac{d^2 i(t)}{dt^2} + R \frac{di(t)}{dt} + \frac{1}{C} i(t) = 0
\end{equation}
The characteristic natural frequency $\omega_0$ and quality factor $Q$ are defined as:
\begin{equation}
\omega_0 = \frac{1}{\sqrt{LC}}, \qquad Q = \frac{1}{R}\sqrt{\frac{L}{C}}
\end{equation}

\section{Apparatus and Experimental Setup}
\begin{itemize}
    \item Digital Storage Oscilloscope (Rigol DS1054Z, $\pm 2\%$ vertical accuracy)
    \item Arbitrary Function Generator (Keysight 33500B)
    \item Precision Decade Inductor Box ($\qty{10.0}{\milli\henry} \pm 1\%$)
    \item Precision Capacitor ($\qty{10.0}{\nano\farad} \pm 1\%$)
    \item Metal Film Resistors ($R_1 = \qty{100.0}{\ohm}$, $R_2 = \qty{470.0}{\ohm}$)
\end{itemize}

\section{Experimental Observations and Data}
Measurements were recorded under a $\qty{5.0}{\volt}$ peak-to-peak square wave excitation:

\begin{table}[htbp]
\centering
\caption{Experimental Transient Response Observations}
\label{tab:measurements}
\begin{tabular}{cccccc}
\toprule
\textbf{Trial} & \textbf{Resistance ($\Omega$)} & \textbf{Peak Voltage (\unit{\volt})} & \textbf{Overshoot (\%)} & \textbf{Period (\unit{\micro\second})} \\
\midrule
1 & $100.2$ & $7.45 \pm 0.05$ & $49.0$ & $64.8 \pm 0.4$ \\
2 & $100.2$ & $7.48 \pm 0.05$ & $49.6$ & $64.7 \pm 0.4$ \\
3 & $100.2$ & $7.42 \pm 0.05$ & $48.4$ & $65.0 \pm 0.4$ \\
4 & $470.5$ & $5.62 \pm 0.05$ & $12.4$ & $68.2 \pm 0.4$ \\
5 & $470.5$ & $5.60 \pm 0.05$ & $12.0$ & $68.4 \pm 0.4$ \\
\bottomrule
\end{tabular}
\end{table}

\section{Uncertainty and Error Propagation}
Given period $T = \qty{64.8 \pm 0.4}{\micro\second}$, the empirical frequency $f$ is:
\begin{equation}
f = \frac{1}{T} = \qty{15.43}{\kilo\hertz}
\end{equation}
Propagating the measurement uncertainty using partial derivatives:
\begin{equation}
\delta f = \left| \frac{\partial f}{\partial T} \right| \delta T = \frac{\delta T}{T^2} = \frac{\qty{0.4}{\micro\second}}{(\qty{64.8}{\micro\second})^2} = \qty{0.095}{\kilo\hertz}
\end{equation}
Accounting for component tolerances ($\delta L/L = 1\%$, $\delta C/C = 1\%$), the theoretical resonance is:
\begin{equation}
f_{theory} = \frac{1}{2\pi \sqrt{LC}} = \qty{15.92 \pm 0.16}{\kilo\hertz}
\end{equation}

\section{Discussion and Sources of Error}
The observed percentage discrepancy of $3.1\%$ is fully encompassed within the cumulative $2\sigma$ error bounds. Systematic errors arose primarily from:
\begin{enumerate}
    \item Oscilloscope probe capacitive loading ($\approx \qty{15}{\pico\farad}$).
    \item Internal ESR of the inductor winding ($\approx \qty{3.2}{\ohm}$).
\end{enumerate}

\section{Conclusion}
The theoretical model for second-order RLC transients was validated experimentally. Damping and frequency measurements were consistent across multiple trials within quantified instrument uncertainty.

\end{document}
